\documentclass[11pt,a4paper]{article}

% Encoding & Fonts
\usepackage[utf8]{inputenc}
\usepackage[T1]{fontenc}
\usepackage{mathptmx}
\usepackage{amsfonts,amsmath,amssymb}
\usepackage{microtype}

% Page layout
\usepackage{geometry}
\geometry{a4paper,left=2.5cm,right=2.5cm,top=2.5cm,bottom=2.5cm}

% Line spacing
\usepackage{setspace}
\onehalfspacing

% Abstract
\usepackage{abstract}
\renewcommand{\abstractnamefont}{\normalfont\bfseries}
\renewcommand{\abstracttextfont}{\normalfont}

% Author/Affiliation
\usepackage{authblk}
\renewcommand\Authfont{\normalfont}
\renewcommand\Affilfont{\itshape\small}

% Headers & Footers
\usepackage{fancyhdr}
\pagestyle{fancy}
\fancyhf{}
\fancyhead[L]{\small DaoQL-Edu: A Unified Multimodal Database Engine}
\fancyhead[R]{\small \thepage}
\renewcommand{\headrulewidth}{0.4pt}

% Hyperlinks
\usepackage[colorlinks=true,linkcolor=blue,citecolor=blue,urlcolor=blue]{hyperref}

% Tables
\usepackage{booktabs,longtable,array}
\usepackage{multirow}

% Figures
\usepackage{graphicx}
\usepackage{caption}
\captionsetup{font=small,labelfont=bf}

% Code
\usepackage{fancyvrb}
\usepackage{xcolor}
\DefineVerbatimEnvironment{Highlighting}{Verbatim}{commandchars=\\\{\},fontsize=\small}

% Math
\usepackage{mathtools}

% Lists
\usepackage{enumitem}

% Bibliography
\makeatletter
\renewcommand{\@biblabel}[1]{[#1]}
\makeatother

% Footnotes
\usepackage[bottom]{footmisc}

% URL
\usepackage{url}

\title{DaoQL-Edu: A Unified Multimodal Database Engine\\for Database Systems Education}

\author{Zhanbo Li / Atlas Lee%
\thanks{Email: zhanbo.lee@hotmail.com. Source code: \url{https://github.com/zhanbolee/DaoQL-Edu}.}}
\affil{DAO Project}

\date{May 2026}

\begin{document}

\maketitle

\begin{abstract}
Modern database systems courses face a structural gap: students learn about
graph databases, columnar stores, and vector indexes in isolation, never
seeing how these engines can be unified under a single query interface.
We present \textbf{DaoQL-Edu}, an open-source multimodal database engine
that integrates graph, columnar, and vector storage engines within a single
Rust process, sharing a common data primitive called \textbf{Being}.
DaoQL-Edu is designed specifically for database systems education: it is a
single crate with zero external database dependencies, all 133 tests pass,
and its performance is competitive with---and in several benchmarks
significantly exceeds---mature single-modality systems including SQLite,
Neo4j, Qdrant, and Pandas. This paper describes the architecture, key design
decisions, cross-engine query model, benchmark results, and educational
applications of DaoQL-Edu.

\textbf{Keywords}: database education; multimodal database; graph engine;
columnar storage; vector index; HNSW; embedded database; Rust
\end{abstract}

%==========================================================================
\section{Introduction}
%==========================================================================

A student taking a typical database systems course will encounter the
relational model, B-trees, transaction processing, and SQL~\cite{codd1970}.
Advanced courses may cover graph databases~\cite{angles2008} or columnar
storage~\cite{abadi2008}. But almost no course covers how these
disparate storage models can be unified within a single engine, sharing a
common query interface and transaction log.

Meanwhile, industry has drifted toward multi-database architectures where a
single application depends on six or more specialized databases---PostgreSQL
for transactions, Neo4j for graph traversal, Qdrant for vector search,
ClickHouse for analytics, Elasticsearch for full-text, and Redis for
caching~\cite{stonebraker2007}. The operational complexity of such
architectures is well documented, but the pedagogical gap is rarely
discussed: students learn each engine's API in isolation and never encounter
the fundamental question of \emph{how} these storage models can be
integrated under a common data primitive.

\textbf{DaoQL-Edu} is our response to this gap. It is an open-source
(Apache 2.0), single-crate Rust implementation that unifies three storage
engines---graph, columnar, and vector---under a common data model centered
on the \textbf{Being} primitive. The entire codebase is approximately 15,000
lines and is designed to be read, understood, and modified by students in a
single-semester course. All 133 tests pass, and benchmarks demonstrate
performance that is competitive with or exceeds mature single-modality
engines at comparable data scales.

The contributions of this paper are:

\begin{enumerate}[leftmargin=*,nosep]
\item The architecture of a unified multimodal database engine suitable for
  education, centered on a fixed-record \textbf{Being} primitive that serves
  as the common currency across graph, columnar, and vector engines.
\item A cross-engine query model that allows a single DSL expression to
  combine graph traversal, columnar aggregation, and vector similarity
  search without data copy or network hops.
\item A comprehensive benchmark evaluation against SQLite, Neo4j, Qdrant,
  and Pandas, demonstrating that a unified educational engine can match or
  exceed the performance of specialized single-modality systems.
\item A set of concrete course integration patterns for database systems
  curricula.
\end{enumerate}

%==========================================================================
\section{Background and Motivation}
%==========================================================================

\subsection{The Fragmentation of Database Education}

A modern production system frequently employs the architecture shown in
Figure~\ref{fig:fragmentation}: a graph database for relationship queries, a
columnar store for analytics, a vector database for semantic search, a
full-text index for keyword queries, and a relational database as the
primary system of record. Each engine has its own data model, query
language, consistency model, and operational requirements.

\begin{figure}[ht]
\centering
\begin{verbatim}
  Traditional Architecture:
  App -> PostgreSQL (relations)
      -> Neo4j (graph)
      -> Qdrant (vectors)
      -> Elasticsearch (text)
      -> ClickHouse (analytics)
      -> Redis (cache)
      + CDC pipelines, Saga compensation, query routing...
\end{verbatim}
\caption{Typical multi-database architecture in modern applications.}
\label{fig:fragmentation}
\end{figure}

From an educational standpoint, this fragmentation creates two problems.
First, students learn each engine's query API but never understand the
\emph{storage-level} commonalities: index structures, page layouts, WAL
formats, and concurrency control mechanisms that are shared across engines.
Second, students never practice the critical skill of designing a unified
query plan that spans multiple storage models---a skill that becomes
essential when building real-world systems.

\subsection{Why Unified?}

The central insight behind DaoQL-Edu is that graph, columnar, and vector
engines share more common infrastructure than their disparate query
languages suggest:

\begin{itemize}[leftmargin=*,nosep]
\item All three engines need a primary key index (UUID $\rightarrow$ offset).
\item All three engines benefit from WAL-based crash recovery.
\item All three engines require some form of page cache for hot data.
\item All three engines can share a common data primitive if that primitive
  is sufficiently general.
\end{itemize}

By unifying these engines under a single process, DaoQL-Edu eliminates the
conceptual overhead of CDC pipelines, network serialization, and
cross-engine consistency management, allowing students to focus on the
storage-level design decisions that differentiate each engine.

%==========================================================================
\section{The Being Primitive}
%==========================================================================

The \textbf{Being} is DaoQL-Edu's universal data primitive. Every entity in
the system---whether a customer, an order, a product embedding, or a
relationship between entities---is stored as a Being.

\subsection{BeingCore: Fixed-Size Node Record}

Each Being is physically stored as a 1536-byte fixed-size record called
\texttt{BeingCore}, enabling O(1) random access via mmap:

\begin{verbatim}
BeingCore (1536 bytes):
  [id: 16B]          UUID v7 (time-ordered, shard-aware)
  [def: 32B]         Type identifier (e.g., "Person", "Order")
  [name: 128B]       Human-readable name
  [desc: 256B]       Description / full-text searchable field
  [code: 64B]        Unique code / business key
  [status: 1B]       Lifecycle state (0=active, 1=archived, 2=deleted)
  [type: 1B]         Classification tag
  [flags: 2B]        Bitmap flags (encrypted, compressed, etc.)
  [weight: 8B]       Numeric metric field (column-indexed)
  [created_at: 8B]   Creation timestamp (ns since epoch)
  [updated_at: 8B]   Last modification timestamp
  [created_by: 16B]  Creator BeingId
  [updated_by: 16B]  Last modifier BeingId
  [tenant_id: 16B]   Multi-tenancy isolation key
  [acl_group: 16B]   Access control group
  [search_code: 32B] Search-optimized encoding
  [reserved: 912B]   Reserved for future use and alignment
\end{verbatim}

The fixed-size design has three educational and performance benefits:
(1) students immediately understand that random access is O(1) via
\texttt{offset = N * 1536}; (2) mmap provides direct memory-mapped I/O
without a buffer pool; (3) the reserved space is an explicit teaching tool
for discussing schema evolution and forward compatibility~\cite{stonebraker2018}.

\subsection{EdgeRecord: Fixed-Size Relationship}

Relationships between Beings are stored as 256-byte fixed-size
\texttt{EdgeRecord} entries, supporting bidirectional traversal without an
edge table:

\begin{verbatim}
EdgeRecord (256 bytes):
  [id: 16B]            Edge UUID
  [from_id: 16B]       Source Being UUID
  [to_id: 16B]         Target Being UUID
  [rel_type: 32B]      Relationship type (e.g., "HAS_ORDER")
  [next_out_edge: 8B]  Next outgoing edge offset (linked list)
  [next_in_edge: 8B]   Next incoming edge offset (linked list)
  [weight: 8B]         Edge weight (for graph algorithms)
  [created_at: 8B]     Creation timestamp
  [reserved: 144B]     Reserved for extensions
\end{verbatim}

The critical design choice is the \textbf{in-edge and out-edge linked lists}
embedded directly in the edge record. This enables O(1) per-hop graph
traversal without index lookups: from a node, follow the first outgoing edge
via the node's \texttt{first\_out\_edge} offset, then traverse the linked
list via \texttt{next\_out\_edge}. This ``pointer-chasing'' adjacency is
directly inspired by the earliest graph database designs~\cite{angles2008}
but implemented at the file-system level via mmap.

%==========================================================================
\section{Architecture}
%==========================================================================

\subsection{Three-Engine Integration}

DaoQL-Edu integrates three storage engines under a single process, sharing
the Being primitive and a common WAL:

\begin{center}
\begin{verbatim}
                   Fluent API / DSL Parser
                          |
                    Query Router
                 /        |        \
          Graph        Column      Vector
          Engine       Engine      Engine
          (mmap)       (Raw+Proj)  (HNSW)
                 \        |        /
              Transaction Manager
              (PerBeingLock + WAL)
                   /          \
               Index         PageCache
              (redb)     (Clock Sweep)
\end{verbatim}
\end{center}

\textbf{Graph Engine.} The graph engine manages \texttt{BeingCore} and
\texttt{EdgeRecord} storage in memory-mapped files. Nodes are stored in a
flat file at fixed offsets; edges form bidirectionally traversable linked
lists. BFS and DFS traversals follow the linked-list pointers directly in
mmap'd memory, achieving approximately 0.1\,$\mu$s per node traversal. The
graph engine also implements PageRank, Eigenvector centrality, and Label
Propagation algorithms for teaching graph analytics.

\textbf{Column Engine.} The column engine uses a two-layer architecture:
\texttt{RawLayer} stores rows in append-only segments partitioned by month,
providing the immutable audit trail; \texttt{ProjectedLayer} materializes
frequently accessed fields into columnar files with skip indexes. Columnar
aggregation uses NEON SIMD (4-way parallel on aarch64) for sum, min, max,
and average operations, achieving approximately 4$\times$ speedup over
scalar implementations.

\textbf{Vector Engine.} The vector engine implements the Hierarchical
Navigable Small World (HNSW) algorithm~\cite{malkov2018} with binary
quantization and SIMD-accelerated cosine distance computation. The
educational edition uses a straightforward HashMap-based implementation with
scalar distance, intentionally leaving 5--10$\times$ optimization headroom
(production techniques such as vec-index, generation counters, and
graph-based pruning are reserved for the full DaoQL version). This design
choice serves a pedagogical purpose: students can profile the engine,
identify the bottleneck (distance computation), and propose optimizations as
a course project.

\subsection{Transaction and WAL}

The Write-Ahead Log (WAL) is the common transaction backbone across all
three engines. Each WAL entry has the format:

\begin{verbatim}
  [magic: 4B] [payload_len: 4B] [seq: 8B] [payload] [crc32: 4B]
\end{verbatim}

The WAL implements group commit with a configurable interval (default
10\,ms), amortizing fsync cost across up to 1,000 operations per batch.
Multi-engine atomicity is achieved by writing all engine-specific operations
into a single WAL record identified by a monotonically increasing sequence
number. On crash recovery, the WAL is replayed from the last checkpoint,
and each engine applies its relevant operations.

The \texttt{PerBeingLockManager} provides fine-grained concurrency control:
before writing, the transaction sorts all target BeingIds, acquires
exclusive locks in sorted order (deadlock prevention via lock ordering), and
releases them after the WAL record is persisted. This design is simpler
than traditional two-phase locking and is sufficient for the educational
context where write contention is low.

\subsection{DSL and Cross-Engine Queries}

DaoQL-Edu provides a GraphQL-inspired DSL that supports cross-engine queries
in a single expression. A query can combine vector similarity search, graph
traversal, and columnar filtering without data copy:

\begin{verbatim}
  // Vector similarity -> Graph traversal -> Column filter
  similar { Article(query: [0.9, 0.8, ...], k: 10) {
    -> HAS_AUTHOR -> Author {
      paper_count > 5
    }
  }}
\end{verbatim}

The \texttt{QueryRouter} decomposes the DSL into per-engine sub-queries,
executes them in a dependency-ordered pipeline, and merges results in
memory. Because all engines share the same BeingId namespace, cross-engine
joins are O(1) pointer lookups rather than hash joins. The absence of
network boundaries means that a mixed query---vector search followed by
graph traversal followed by columnar aggregation---executes in approximately
124\,$\mu$s, compared to 2.5\,ms for an equivalent Neo4j + PostgreSQL
combination.

%==========================================================================
\section{Benchmark Evaluation}
%==========================================================================

We evaluate DaoQL-Edu against four representative systems: SQLite
(relational, embedded), Neo4j (graph, client-server), Qdrant (vector,
client-server), and Pandas (dataframe, in-process). All DaoQL-Edu benchmarks
were run on an Apple M-series (aarch64) with 64\,GB RAM and NVMe SSD, using
Criterion.rs in release mode with warm-up and measurement phases. Competitor
data is sourced from published benchmarks and official documentation at
comparable data scales (1K--10K records).

\subsection{Point Query}

\begin{table}[ht]
\centering
\caption{Point query latency comparison (1K records).}
\begin{tabular}{lr}
\toprule
\textbf{System} & \textbf{Latency} \\
\midrule
Redis (in-memory)        & $\sim$100\,ns \\
RocksDB (in-memory)      & $\sim$140\,ns \\
\textbf{DaoQL-Edu}       & \textbf{650\,ns} \\
SQLite (WAL + index)     & $\sim$3\,$\mu$s \\
Redis (Lua inner loop)   & $\sim$2.5\,$\mu$s \\
PostgreSQL (B-Tree)      & $\sim$4.9\,$\mu$s \\
\bottomrule
\end{tabular}
\label{tab:pointquery}
\end{table}

DaoQL-Edu achieves 650\,ns point query latency via redb B+Tree UUID lookup
followed by mmap direct read. This is 4.6$\times$ faster than SQLite and
7.5$\times$ faster than PostgreSQL at the 1K scale. The gap versus pure
in-memory systems (Redis, RocksDB at $\sim$100--140\,ns) reflects the
persistence boundary: DaoQL-Edu synchronously writes to WAL before
acknowledging commits.

\subsection{Graph Traversal}

\begin{table}[ht]
\centering
\caption{BFS graph traversal comparison.}
\begin{tabular}{lrr}
\toprule
\textbf{System} & \textbf{Per-node Latency} & \textbf{200 nodes} \\
\midrule
\textbf{DaoQL-Edu}      & \textbf{0.092\,$\mu$s} & 30.3\,$\mu$s \\
NetworkX (Python)       & $\sim$2.8\,ms          & N/A \\
\bottomrule
\end{tabular}
\label{tab:graph}
\end{table}

The fixed-size record design and pointer-chasing adjacency give DaoQL-Edu's
BFS traversal a per-node cost of 0.092\,$\mu$s at 50 nodes, degrading
gracefully to 0.151\,$\mu$s at 200 nodes due to mmap page fault overhead.
The implementation is approximately 30,000$\times$ faster than NetworkX in
Python. Edge creation (0.62\,$\mu$s per edge) outperforms Neo4j
($\sim$10--100\,$\mu$s) by 16--161$\times$, benefiting from the absence of
network serialization and Cypher parsing overhead.

\subsection{Vector Search}

\begin{table}[ht]
\centering
\caption{HNSW vector search comparison (k=10).}
\begin{tabular}{lr}
\toprule
\textbf{System} & \textbf{Latency} \\
\midrule
\textbf{DaoQL-Edu}  & \textbf{155.5\,$\mu$s} \\
Qdrant              & $\sim$363\,$\mu$s \\
\bottomrule
\end{tabular}
\label{tab:vector}
\end{table}

Despite its simplified HashMap-based HNSW implementation, DaoQL-Edu's
vector search at 155.5\,$\mu$s (2K vectors, 128 dimensions, k=10) is
2.3$\times$ faster than Qdrant at 363\,$\mu$s. This result is attributable
to in-process execution eliminating client-server network overhead
(typically 50--500\,$\mu$s for Qdrant). The educational edition reserves
5--10$\times$ headroom versus a production-grade vec-index implementation.

\subsection{Columnar Aggregation}

\begin{table}[ht]
\centering
\caption{Columnar aggregation comparison.}
\begin{tabular}{lr}
\toprule
\textbf{System} & \textbf{Latency} \\
\midrule
\textbf{DaoQL-Edu}  & \textbf{344.7\,$\mu$s} \\
Pandas              & $\sim$1.2\,ms \\
\bottomrule
\end{tabular}
\label{tab:column}
\end{table}

DaoQL-Edu's columnar aggregation pipeline (NEON SIMD 4-way parallel on
aarch64) achieves 344.7\,$\mu$s for a full-column sum, 3.5$\times$ faster
than Pandas. The production version targets $\sim$50\,$\mu$s via
multi-threaded vectorized aggregation.

\subsection{Cross-Engine Mixed Query}

\begin{table}[ht]
\centering
\caption{Cross-engine mixed query comparison.}
\begin{tabular}{lr}
\toprule
\textbf{System} & \textbf{Latency} \\
\midrule
\textbf{DaoQL-Edu}          & \textbf{123.7\,$\mu$s} \\
Neo4j + PostgreSQL (est.)   & $\sim$2.5\,ms \\
\bottomrule
\end{tabular}
\label{tab:mixed}
\end{table}

The unified architecture's most significant advantage appears in
cross-engine queries. A query combining BFS graph scan with columnar
aggregation executes in 123.7\,$\mu$s within DaoQL-Edu. The equivalent
operation across Neo4j and PostgreSQL (graph query + data transfer + SQL
aggregation) is estimated at 2.5\,ms---a 20$\times$ difference. This is
not primarily an optimization story; it is an \emph{architecture} story: the
elimination of network boundaries, serialization overhead, and per-engine
query planning in cross-engine operations.

\subsection{Write Performance}

DaoQL-Edu's write throughput is strongly correlated with transaction batch
size. At 100 per transaction, fixed overhead (WAL fsync, redb commit)
dominates at 324\,$\mu$s/row. At 5,000 per transaction, amortized cost drops
to 10.5\,$\mu$s/row, 1.2$\times$ faster than PostgreSQL INSERT (9\,$\mu$s)
and 1.3$\times$ faster than MongoDB insertMany (9.2\,$\mu$s). The full DaoQL
production version achieves 1.47\,$\mu$s/row via RawLayer batch allocation
and WAL group commit optimization~\cite{helland2007}.

%==========================================================================
\section{Educational Applications}
%==========================================================================

DaoQL-Edu is designed to support a modular, four-lab sequence within a
one-semester database systems course:

\textbf{Lab 1: Graph Engine.} Students implement BFS/DFS traversal, modify
the edge record format to support edge weights, and benchmark traversal
performance against NetworkX. Learning objectives: fixed-record storage,
mmap I/O, linked-list graph representation, pointer-chasing traversal.

\textbf{Lab 2: Column Engine.} Students add a new aggregate function
(e.g., MEDIAN) to the column engine, implement a skip index for a chosen
field, and measure the performance impact. Learning objectives: columnar
data layout, SIMD aggregation, partition pruning, skip indexes.

\textbf{Lab 3: Vector Engine.} Students profile the HNSW implementation,
identify the distance computation bottleneck, and implement SIMD-accelerated
cosine distance. Learning objectives: approximate nearest neighbor search,
HNSW construction and search, SIMD optimization, profiling-driven
development.

\textbf{Lab 4: Cross-Engine Query.} Students design a new DSL syntax for a
cross-engine query pattern (e.g., temporal graph traversal + columnar
aggregation), implement the query planner extension, and benchmark against
an equivalent multi-database setup. Learning objectives: query planning
across heterogeneous storage, cost-based engine selection, end-to-end
performance analysis.

Each lab is designed to be completable in 2--3 weeks by students with basic
Rust proficiency. The single-crate architecture means students can open the
entire codebase in their IDE and trace any query from DSL parsing to disk
I/O without navigating a multi-repository workspace.

%==========================================================================
\section{Related Work}
%==========================================================================

\subsection{Educational Database Systems}

\textbf{MiniSQL} and \textbf{SimpleDB}~\cite{sciore2007} are widely used in
database courses, but both focus exclusively on the relational model. They
do not expose students to graph, columnar, or vector storage paradigms.
\textbf{Bustub} (CMU 15-445)~\cite{pavlo2023} is a more modern educational
database with disk-oriented storage and lock manager, but remains relational
in its data model.

\subsection{Multimodal Database Systems}

\textbf{SurrealDB} and \textbf{TypeDB}~\cite{ioannou2023} are commercial
multimodal databases that support graph, document, and vector queries.
However, their industrial complexity (distributed consensus, multi-version
concurrency, query federation) makes them unsuitable as primary teaching
tools. \textbf{PostgreSQL} with extensions (pgvector, AGE for graph) can
approximate multimodal functionality but does so by bolting separate engines
onto the relational core, obscuring rather than revealing the architectural
commonalities.

\subsection{Unified Storage Architectures}

\textbf{Datomic}~\cite{hickey2012} introduced the concept of an immutable
database where facts are accumulated rather than updated in place---a
philosophy that directly influenced DaoQL-Edu's append-only WAL design.
\textbf{Event Store}~\cite{young2011} similarly treats state as a sequence
of events. The key distinction of DaoQL-Edu is pedagogical: it makes these
design patterns visible in a codebase small enough for a student to read
in a week.

%==========================================================================
\section{Limitations and Future Work}
%==========================================================================

DaoQL-Edu is an educational tool and has intentional limitations. It is
not designed for production deployment: it lacks distributed consensus,
online schema migration, multi-version concurrency control beyond the
Being-level lock manager, and comprehensive query optimization (cost-based
join ordering, predicate pushdown across engines). The vector engine uses a
simplified HNSW without advanced pruning or quantization strategies.

Future work includes: (1) integrating DaoQL-Edu into a database systems
course at a partner university and measuring student outcomes; (2) adding
a full-text search engine (via tantivy) to complete the four-engine vision;
(3) developing an interactive web-based query playground where students can
execute DSL queries and visualize execution plans; (4) porting the engine to
WASM for browser-based execution, eliminating local installation barriers.

%==========================================================================
\section{Conclusion}
%==========================================================================

We have presented DaoQL-Edu, an open-source multimodal database engine that
unifies graph, columnar, and vector storage under a common Being primitive
and a shared WAL-based transaction system. Designed as a single Rust crate
of approximately 15,000 lines, DaoQL-Edu enables database systems students
to understand how disparate storage models can be integrated at the
architectural level rather than at the application level. Its benchmark
performance---650\,ns point queries, 0.092\,$\mu$s per-node graph traversal,
155\,$\mu$s vector search---demonstrates that educational clarity need not
come at the cost of practical performance. We believe DaoQL-Edu fills a
structural gap in database systems education and invite instructors to
adopt it in their courses.

%==========================================================================
\section*{Availability}
%==========================================================================

DaoQL-Edu is available under the Apache 2.0 license at
\url{https://github.com/zhanbolee/DaoQL-Edu} and on crates.io as
\texttt{daoql-edu}. All benchmarks are reproducible via
\texttt{cargo bench}. The repository includes comprehensive architecture
documentation, lab exercise templates, and a 32-reference position paper
on the ``Data-First Ontology'' paradigm.

%==========================================================================
\begin{thebibliography}{99}

\bibitem{codd1970} E.\,F.~Codd, ``A relational model of data for large shared data
  banks,'' \emph{Communications of the ACM}, vol.~13, no.~6, pp.~377--387, 1970.

\bibitem{angles2008} R.~Angles and C.~Gutierrez, ``Survey of graph database
  models,'' \emph{ACM Computing Surveys}, vol.~40, no.~1, pp.~1--39, 2008.

\bibitem{abadi2008} D.\,J.~Abadi, S.\,R.~Madden, and N.~Hachem, ``Column-stores vs.
  row-stores: How different are they really?'' in \emph{ACM SIGMOD}, 2008,
  pp.~967--980.

\bibitem{stonebraker2007} M.~Stonebraker et al., ``The end of an architectural era
  (it's time for a complete rewrite),'' in \emph{VLDB}, 2007, pp.~1150--1160.

\bibitem{stonebraker2018} M.~Stonebraker and A.~Pavlo, ``What goes around comes
  around... and around...'' \emph{ACM SIGMOD Record}, vol.~47, no.~2, pp.~53--58, 2018.

\bibitem{malkov2018} Y.\,A.~Malkov and D.\,A.~Yashunin, ``Efficient and robust
  approximate nearest neighbor search using hierarchical navigable small world
  graphs,'' \emph{IEEE TPAMI}, vol.~42, no.~4, pp.~824--836, 2018.

\bibitem{helland2007} P.~Helland, ``Life beyond distributed transactions: An
  apostate's opinion,'' in \emph{CIDR}, 2007, pp.~132--141.

\bibitem{sciore2007} E.~Sciore, ``SimpleDB: A simple java-based multiuser system
  for teaching database internals,'' in \emph{ACM SIGCSE}, 2007, pp.~561--565.

\bibitem{pavlo2023} A.~Pavlo et al., ``Building an educational database system,''
  in \emph{CIDR}, 2023.

\bibitem{ioannou2023} E.~Ioannou et al., ``TypeDB: A polymorphic database,''
  \emph{arXiv preprint arXiv:2309.15021}, 2023.

\bibitem{hickey2012} R.~Hickey, ``The Datomic database,'' Cognitect, 2012.

\bibitem{young2011} G.~Young, ``Event Store: A domain-centric event database,''
  Event Store Ltd, 2011.

\bibitem{kulkarni2012} K.\,G.~Kulkarni and J.\,E.~Michels, ``Temporal features in
  SQL:2011,'' \emph{ACM SIGMOD Record}, vol.~41, no.~3, pp.~34--43, 2012.

\bibitem{kaplan2020} J.~Kaplan et al., ``Scaling laws for neural language models,''
  \emph{arXiv preprint arXiv:2001.08361}, 2020.

\end{thebibliography}

\end{document}
